\documentclass{article}
\usepackage[utf8]{inputenc}
\usepackage[spanish]{babel}
\usepackage{booktabs} % INDISPENSABLE para tablas de Pandas
\usepackage{graphicx}

% Graficos de Matplot
\newcommand{\mathdefault}[1][]{#1}
\usepackage{pgfplots}
\usepackage{tikz}
\usepgfplotslibrary{fillbetween}
\usetikzlibrary{positioning, arrows.meta}
\pgfplotsset{compat=1.18}
\usepackage{import}

\title{Práctica de Estadística Descriptiva}
\author{Gino Anci}

\begin{document}
\maketitle

\section{Tabla de frecuencias}
A continuación se presentan los estadísticos descriptivos obtenidos de los empleados de una determinada empresa:

\begin{table}[h]
\centering
\caption{Resumen Estadístico de la Muestra}
\label{tab:df_tabla_frecuencias}
\begin{tabular}{cccccccc}
\toprule
 & Li & Ls & ni & ci & fi & Fi & Ni \\
\midrule
0 & 800 & 1000 & 51 & 900.000000 & 0.237209 & 0.237209 & 51 \\
1 & 1000 & 1200 & 87 & 1100.000000 & 0.404651 & 0.641860 & 138 \\
2 & 1200 & 1400 & 37 & 1300.000000 & 0.172093 & 0.813953 & 175 \\
3 & 1400 & 1600 & 22 & 1500.000000 & 0.102326 & 0.916279 & 197 \\
4 & 1600 & 1800 & 12 & 1700.000000 & 0.055814 & 0.972093 & 209 \\
5 & 1800 & 2000 & 6 & 1900.000000 & 0.027907 & 1.000000 & 215 \\
\bottomrule
\end{tabular}
\end{table}


En la tabla \ref{tab:df_tabla_frecuencias}:
\begin{itemize}
    \item \texttt{Li}: representa el limite inferior del salario.
    \item \texttt{Ls}: representa el limite superior del salario.
    \item \texttt{ni}: representa la cantidad de empleados.
    \item \texttt{ci}: representa el punto medio de intervalos.
    \item \texttt{fi}: representa la frecuencia absoluta.
    \item \texttt{Fi}: representa la frecuencia acumulada.
    \item \texttt{Ni}: representa la cantidad acumulada de empleados.
\end{itemize}

\subsection{Medidas de centralización}

Tras procesar los datos de frecuencia, se obtuvieron los siguientes estadísticos:
\newcommand{\media}{1183.72}
\newcommand{\desviacion}{252.37}
\newcommand{\varianza}{63688.48}


\begin{itemize}
    \item \textit{Media aritmética ($\bar{x}$):} \media
    \item \textit{Varianza ($\sigma^2$):} \varianza
    \item \textit{Desviación estándar ($\sigma$):} \desviacion
\end{itemize}

\section{Visualización}
\begin{figure}[!ht]
    \caption{Distribucion de salarios en base a cantidad de empleados}
    \label{fig:hist}
    \input{graficos/distribucion.pgf}
\end{figure}

\section{Medidas}

\subsection{Respecto a los impuestos}

Una medida gubernamental determina que a partir del próximo mes se retendrán impuestos a los salarios que superen los \mathdollar{1500} mensuales. 

\emph{¿Qué porcentaje de los empleados se verán afectados por dicha retención?}

\newcommand{\Proporción}{0.13}
\newcommand{\Porcentaje}{13.49}

La cantidad de empleados que tendran una retención por impuestos superiores a \mathdollar{1500} es de \(\Proporción = \Porcentaje\%\). Ver figura\ref{fig:torta_impuesto_1500} 
\begin{figure}[!ht]
    \label{fig:torta_impuesto_1500}
    %% Creator: Matplotlib, PGF backend
%%
%% To include the figure in your LaTeX document, write
%%   \input{<filename>.pgf}
%%
%% Make sure the required packages are loaded in your preamble
%%   \usepackage{pgf}
%%
%% Also ensure that all the required font packages are loaded; for instance,
%% the lmodern package is sometimes necessary when using math font.
%%   \usepackage{lmodern}
%%
%% Figures using additional raster images can only be included by \input if
%% they are in the same directory as the main LaTeX file. For loading figures
%% from other directories you can use the `import` package
%%   \usepackage{import}
%%
%% and then include the figures with
%%   \import{<path to file>}{<filename>.pgf}
%%
%% Matplotlib used the following preamble
%%   \def\mathdefault#1{#1}
%%   \everymath=\expandafter{\the\everymath\displaystyle}
%%   \IfFileExists{scrextend.sty}{
%%     \usepackage[fontsize=10.000000pt]{scrextend}
%%   }{
%%     \renewcommand{\normalsize}{\fontsize{10.000000}{12.000000}\selectfont}
%%     \normalsize
%%   }
%%   \usepackage{fontspec}
%%   \setmainfont{CMU Serif}
%%   \makeatletter\@ifpackageloaded{underscore}{}{\usepackage[strings]{underscore}}\makeatother
%%
\begingroup%
\makeatletter%
\begin{pgfpicture}%
\pgfpathrectangle{\pgfpointorigin}{\pgfqpoint{5.057065in}{2.709000in}}%
\pgfusepath{use as bounding box, clip}%
\begin{pgfscope}%
\pgfsetbuttcap%
\pgfsetmiterjoin%
\definecolor{currentfill}{rgb}{1.000000,1.000000,1.000000}%
\pgfsetfillcolor{currentfill}%
\pgfsetlinewidth{0.000000pt}%
\definecolor{currentstroke}{rgb}{1.000000,1.000000,1.000000}%
\pgfsetstrokecolor{currentstroke}%
\pgfsetdash{}{0pt}%
\pgfpathmoveto{\pgfqpoint{-0.000000in}{0.000000in}}%
\pgfpathlineto{\pgfqpoint{5.057065in}{0.000000in}}%
\pgfpathlineto{\pgfqpoint{5.057065in}{2.709000in}}%
\pgfpathlineto{\pgfqpoint{-0.000000in}{2.709000in}}%
\pgfpathlineto{\pgfqpoint{-0.000000in}{0.000000in}}%
\pgfpathclose%
\pgfusepath{fill}%
\end{pgfscope}%
\begin{pgfscope}%
\pgfsetbuttcap%
\pgfsetmiterjoin%
\definecolor{currentfill}{rgb}{0.036471,0.140000,0.211765}%
\pgfsetfillcolor{currentfill}%
\pgfsetfillopacity{0.500000}%
\pgfsetlinewidth{1.505625pt}%
\definecolor{currentstroke}{rgb}{0.036471,0.140000,0.211765}%
\pgfsetstrokecolor{currentstroke}%
\pgfsetstrokeopacity{0.500000}%
\pgfsetdash{}{0pt}%
\pgfpathmoveto{\pgfqpoint{3.546673in}{1.248330in}}%
\pgfpathcurveto{\pgfqpoint{3.546673in}{1.379331in}}{\pgfqpoint{3.518813in}{1.508847in}}{\pgfqpoint{3.464948in}{1.628262in}}%
\pgfpathcurveto{\pgfqpoint{3.411083in}{1.747676in}}{\pgfqpoint{3.332433in}{1.854282in}}{\pgfqpoint{3.234231in}{1.940986in}}%
\pgfpathlineto{\pgfqpoint{2.622673in}{1.248330in}}%
\pgfpathlineto{\pgfqpoint{3.546673in}{1.248330in}}%
\pgfpathclose%
\pgfusepath{stroke,fill}%
\end{pgfscope}%
\begin{pgfscope}%
\pgfsetbuttcap%
\pgfsetmiterjoin%
\definecolor{currentfill}{rgb}{0.300000,0.149412,0.016471}%
\pgfsetfillcolor{currentfill}%
\pgfsetfillopacity{0.500000}%
\pgfsetlinewidth{1.505625pt}%
\definecolor{currentstroke}{rgb}{0.300000,0.149412,0.016471}%
\pgfsetstrokecolor{currentstroke}%
\pgfsetstrokeopacity{0.500000}%
\pgfsetdash{}{0pt}%
\pgfpathmoveto{\pgfqpoint{3.192117in}{1.921989in}}%
\pgfpathcurveto{\pgfqpoint{3.113490in}{1.991410in}}{\pgfqpoint{3.023614in}{2.046939in}}{\pgfqpoint{2.926347in}{2.086192in}}%
\pgfpathcurveto{\pgfqpoint{2.829081in}{2.125444in}}{\pgfqpoint{2.725838in}{2.147849in}}{\pgfqpoint{2.621050in}{2.152446in}}%
\pgfpathcurveto{\pgfqpoint{2.516263in}{2.157042in}}{\pgfqpoint{2.411454in}{2.143763in}}{\pgfqpoint{2.311125in}{2.113178in}}%
\pgfpathcurveto{\pgfqpoint{2.210795in}{2.082593in}}{\pgfqpoint{2.116402in}{2.035147in}}{\pgfqpoint{2.031998in}{1.972877in}}%
\pgfpathcurveto{\pgfqpoint{1.947595in}{1.910607in}}{\pgfqpoint{1.874407in}{1.834418in}}{\pgfqpoint{1.815578in}{1.747581in}}%
\pgfpathcurveto{\pgfqpoint{1.756749in}{1.660744in}}{\pgfqpoint{1.713133in}{1.564521in}}{\pgfqpoint{1.686604in}{1.463044in}}%
\pgfpathcurveto{\pgfqpoint{1.660074in}{1.361566in}}{\pgfqpoint{1.651016in}{1.256309in}}{\pgfqpoint{1.659818in}{1.151791in}}%
\pgfpathcurveto{\pgfqpoint{1.668620in}{1.047273in}}{\pgfqpoint{1.695155in}{0.945013in}}{\pgfqpoint{1.738283in}{0.849402in}}%
\pgfpathcurveto{\pgfqpoint{1.781411in}{0.753791in}}{\pgfqpoint{1.840506in}{0.666218in}}{\pgfqpoint{1.913030in}{0.590443in}}%
\pgfpathcurveto{\pgfqpoint{1.985553in}{0.514668in}}{\pgfqpoint{2.070453in}{0.451793in}}{\pgfqpoint{2.164082in}{0.404516in}}%
\pgfpathcurveto{\pgfqpoint{2.257711in}{0.357240in}}{\pgfqpoint{2.358710in}{0.326249in}}{\pgfqpoint{2.462742in}{0.312875in}}%
\pgfpathcurveto{\pgfqpoint{2.566774in}{0.299501in}}{\pgfqpoint{2.672327in}{0.303938in}}{\pgfqpoint{2.774870in}{0.325995in}}%
\pgfpathcurveto{\pgfqpoint{2.877412in}{0.348053in}}{\pgfqpoint{2.975454in}{0.387410in}}{\pgfqpoint{3.064786in}{0.442377in}}%
\pgfpathcurveto{\pgfqpoint{3.154117in}{0.497344in}}{\pgfqpoint{3.233441in}{0.567123in}}{\pgfqpoint{3.299349in}{0.648716in}}%
\pgfpathcurveto{\pgfqpoint{3.365258in}{0.730310in}}{\pgfqpoint{3.416795in}{0.822533in}}{\pgfqpoint{3.451747in}{0.921426in}}%
\pgfpathcurveto{\pgfqpoint{3.486699in}{1.020319in}}{\pgfqpoint{3.504559in}{1.124446in}}{\pgfqpoint{3.504559in}{1.229333in}}%
\pgfpathlineto{\pgfqpoint{2.580559in}{1.229333in}}%
\pgfpathlineto{\pgfqpoint{3.192117in}{1.921989in}}%
\pgfpathclose%
\pgfusepath{stroke,fill}%
\end{pgfscope}%
\begin{pgfscope}%
\pgfsetbuttcap%
\pgfsetmiterjoin%
\definecolor{currentfill}{rgb}{0.121569,0.466667,0.705882}%
\pgfsetfillcolor{currentfill}%
\pgfsetlinewidth{1.505625pt}%
\definecolor{currentstroke}{rgb}{1.000000,1.000000,1.000000}%
\pgfsetstrokecolor{currentstroke}%
\pgfsetdash{}{0pt}%
\pgfpathmoveto{\pgfqpoint{3.572339in}{1.273997in}}%
\pgfpathcurveto{\pgfqpoint{3.572339in}{1.404998in}}{\pgfqpoint{3.544480in}{1.534514in}}{\pgfqpoint{3.490615in}{1.653928in}}%
\pgfpathcurveto{\pgfqpoint{3.436749in}{1.773343in}}{\pgfqpoint{3.358100in}{1.879948in}}{\pgfqpoint{3.259898in}{1.966653in}}%
\pgfpathlineto{\pgfqpoint{2.648339in}{1.273997in}}%
\pgfpathlineto{\pgfqpoint{3.572339in}{1.273997in}}%
\pgfpathclose%
\pgfusepath{stroke,fill}%
\end{pgfscope}%
\begin{pgfscope}%
\pgfsetbuttcap%
\pgfsetmiterjoin%
\definecolor{currentfill}{rgb}{1.000000,0.498039,0.054902}%
\pgfsetfillcolor{currentfill}%
\pgfsetlinewidth{1.505625pt}%
\definecolor{currentstroke}{rgb}{1.000000,1.000000,1.000000}%
\pgfsetstrokecolor{currentstroke}%
\pgfsetdash{}{0pt}%
\pgfpathmoveto{\pgfqpoint{3.217784in}{1.947656in}}%
\pgfpathcurveto{\pgfqpoint{3.139157in}{2.017077in}}{\pgfqpoint{3.049280in}{2.072606in}}{\pgfqpoint{2.952014in}{2.111858in}}%
\pgfpathcurveto{\pgfqpoint{2.854748in}{2.151111in}}{\pgfqpoint{2.751504in}{2.173516in}}{\pgfqpoint{2.646717in}{2.178112in}}%
\pgfpathcurveto{\pgfqpoint{2.541930in}{2.182709in}}{\pgfqpoint{2.437121in}{2.169429in}}{\pgfqpoint{2.336791in}{2.138845in}}%
\pgfpathcurveto{\pgfqpoint{2.236462in}{2.108260in}}{\pgfqpoint{2.142068in}{2.060814in}}{\pgfqpoint{2.057665in}{1.998544in}}%
\pgfpathcurveto{\pgfqpoint{1.973261in}{1.936274in}}{\pgfqpoint{1.900074in}{1.860085in}}{\pgfqpoint{1.841245in}{1.773248in}}%
\pgfpathcurveto{\pgfqpoint{1.782416in}{1.686411in}}{\pgfqpoint{1.738800in}{1.590188in}}{\pgfqpoint{1.712270in}{1.488710in}}%
\pgfpathcurveto{\pgfqpoint{1.685741in}{1.387233in}}{\pgfqpoint{1.676683in}{1.281975in}}{\pgfqpoint{1.685485in}{1.177457in}}%
\pgfpathcurveto{\pgfqpoint{1.694287in}{1.072939in}}{\pgfqpoint{1.720822in}{0.970679in}}{\pgfqpoint{1.763950in}{0.875068in}}%
\pgfpathcurveto{\pgfqpoint{1.807078in}{0.779457in}}{\pgfqpoint{1.866173in}{0.691884in}}{\pgfqpoint{1.938696in}{0.616110in}}%
\pgfpathcurveto{\pgfqpoint{2.011220in}{0.540335in}}{\pgfqpoint{2.096120in}{0.477459in}}{\pgfqpoint{2.189749in}{0.430183in}}%
\pgfpathcurveto{\pgfqpoint{2.283378in}{0.382907in}}{\pgfqpoint{2.384377in}{0.351916in}}{\pgfqpoint{2.488409in}{0.338542in}}%
\pgfpathcurveto{\pgfqpoint{2.592441in}{0.325168in}}{\pgfqpoint{2.697994in}{0.329605in}}{\pgfqpoint{2.800536in}{0.351662in}}%
\pgfpathcurveto{\pgfqpoint{2.903079in}{0.373719in}}{\pgfqpoint{3.001121in}{0.413077in}}{\pgfqpoint{3.090452in}{0.468044in}}%
\pgfpathcurveto{\pgfqpoint{3.179784in}{0.523011in}}{\pgfqpoint{3.259107in}{0.592790in}}{\pgfqpoint{3.325016in}{0.674383in}}%
\pgfpathcurveto{\pgfqpoint{3.390925in}{0.755977in}}{\pgfqpoint{3.442462in}{0.848200in}}{\pgfqpoint{3.477414in}{0.947093in}}%
\pgfpathcurveto{\pgfqpoint{3.512366in}{1.045986in}}{\pgfqpoint{3.530225in}{1.150112in}}{\pgfqpoint{3.530225in}{1.255000in}}%
\pgfpathlineto{\pgfqpoint{2.606225in}{1.255000in}}%
\pgfpathlineto{\pgfqpoint{3.217784in}{1.947656in}}%
\pgfpathclose%
\pgfusepath{stroke,fill}%
\end{pgfscope}%
\begin{pgfscope}%
\definecolor{textcolor}{rgb}{0.000000,0.000000,0.000000}%
\pgfsetstrokecolor{textcolor}%
\pgfsetfillcolor{textcolor}%
\pgftext[x=3.574842in,y=1.691921in,left,]{\color{textcolor}{\rmfamily\fontsize{10.000000}{12.000000}\selectfont\catcode`\^=\active\def^{\ifmmode\sp\else\^{}\fi}\catcode`\%=\active\def%{\%}Sometidos a impuestos}}%
\end{pgfscope}%
\begin{pgfscope}%
\definecolor{textcolor}{rgb}{0.000000,0.000000,0.000000}%
\pgfsetstrokecolor{textcolor}%
\pgfsetfillcolor{textcolor}%
\pgftext[x=3.153704in,y=1.501956in,,]{\color{textcolor}{\rmfamily\fontsize{10.000000}{12.000000}\selectfont\catcode`\^=\active\def^{\ifmmode\sp\else\^{}\fi}\catcode`\%=\active\def%{\%}13.5\%}}%
\end{pgfscope}%
\begin{pgfscope}%
\definecolor{textcolor}{rgb}{0.000000,0.000000,0.000000}%
\pgfsetstrokecolor{textcolor}%
\pgfsetfillcolor{textcolor}%
\pgftext[x=1.679723in,y=0.837075in,right,]{\color{textcolor}{\rmfamily\fontsize{10.000000}{12.000000}\selectfont\catcode`\^=\active\def^{\ifmmode\sp\else\^{}\fi}\catcode`\%=\active\def%{\%}No sometidos a impuestos}}%
\end{pgfscope}%
\begin{pgfscope}%
\definecolor{textcolor}{rgb}{0.000000,0.000000,0.000000}%
\pgfsetstrokecolor{textcolor}%
\pgfsetfillcolor{textcolor}%
\pgftext[x=2.100860in,y=1.027041in,,]{\color{textcolor}{\rmfamily\fontsize{10.000000}{12.000000}\selectfont\catcode`\^=\active\def^{\ifmmode\sp\else\^{}\fi}\catcode`\%=\active\def%{\%}86.5\%}}%
\end{pgfscope}%
\begin{pgfscope}%
\definecolor{textcolor}{rgb}{0.000000,0.000000,0.000000}%
\pgfsetstrokecolor{textcolor}%
\pgfsetfillcolor{textcolor}%
\pgftext[x=2.606225in,y=2.493333in,,base]{\color{textcolor}{\rmfamily\fontsize{12.000000}{14.400000}\selectfont\catcode`\^=\active\def^{\ifmmode\sp\else\^{}\fi}\catcode`\%=\active\def%{\%}Empleados con respecto al impuesto 1500}}%
\end{pgfscope}%
\end{pgfpicture}%
\makeatother%
\endgroup%

\end{figure}

\subsection{Aumento al 12\% de los empleados}

Tras reclamos continuos del personal y ante la imposibilidad de conceder un aumento general, la empresa decidio otorgar un aumento de \mathdollar{120} mensuales a los empleados que menos ganan. 
La empresa lo va a otorgar al 12\% con menores ingresos \emph{¿Cuánto es el salario máximo para obtener dicho aumento?}

\newcommand{\InterpolacionPdoce}{901.18}
\newcommand{\InterpolacionPveinticinco}{1006.32}
\newcommand{\PorcentajeMenorAMil}{23.72}


El salario máximo para obtener dicho aumento es {\InterpolacionPdoce}.

\subsubsection{Caso arreglo con sindicato}

Hay una probabilidad que el sindicato negocie con la empresa y consiga que el aumento se otorge a una cuarta parte de los empleados.

\emph{¿Cuánto es el salario máximo en ese caso?}

Se determino que el salario máximo en caso de arreglo con sindicato es \InterpolacionPveinticinco

Ver comparación del 12\% contra el 25\% en \ref{fig:aumentosalarial}

\begin{figure}[!ht]
    \label{fig:aumentosalarial}
    %% Creator: Matplotlib, PGF backend
%%
%% To include the figure in your LaTeX document, write
%%   \input{<filename>.pgf}
%%
%% Make sure the required packages are loaded in your preamble
%%   \usepackage{pgf}
%%
%% Also ensure that all the required font packages are loaded; for instance,
%% the lmodern package is sometimes necessary when using math font.
%%   \usepackage{lmodern}
%%
%% Figures using additional raster images can only be included by \input if
%% they are in the same directory as the main LaTeX file. For loading figures
%% from other directories you can use the `import` package
%%   \usepackage{import}
%%
%% and then include the figures with
%%   \import{<path to file>}{<filename>.pgf}
%%
%% Matplotlib used the following preamble
%%   \def\mathdefault#1{#1}
%%   \everymath=\expandafter{\the\everymath\displaystyle}
%%   \IfFileExists{scrextend.sty}{
%%     \usepackage[fontsize=10.000000pt]{scrextend}
%%   }{
%%     \renewcommand{\normalsize}{\fontsize{10.000000}{12.000000}\selectfont}
%%     \normalsize
%%   }
%%   \usepackage{fontspec}
%%   \setmainfont{CMU Serif}
%%   \makeatletter\@ifpackageloaded{underscore}{}{\usepackage[strings]{underscore}}\makeatother
%%
\begingroup%
\makeatletter%
\begin{pgfpicture}%
\pgfpathrectangle{\pgfpointorigin}{\pgfqpoint{3.853890in}{3.701083in}}%
\pgfusepath{use as bounding box, clip}%
\begin{pgfscope}%
\pgfsetbuttcap%
\pgfsetmiterjoin%
\definecolor{currentfill}{rgb}{1.000000,1.000000,1.000000}%
\pgfsetfillcolor{currentfill}%
\pgfsetlinewidth{0.000000pt}%
\definecolor{currentstroke}{rgb}{1.000000,1.000000,1.000000}%
\pgfsetstrokecolor{currentstroke}%
\pgfsetdash{}{0pt}%
\pgfpathmoveto{\pgfqpoint{0.000000in}{0.000000in}}%
\pgfpathlineto{\pgfqpoint{3.853890in}{0.000000in}}%
\pgfpathlineto{\pgfqpoint{3.853890in}{3.701083in}}%
\pgfpathlineto{\pgfqpoint{0.000000in}{3.701083in}}%
\pgfpathlineto{\pgfqpoint{0.000000in}{0.000000in}}%
\pgfpathclose%
\pgfusepath{fill}%
\end{pgfscope}%
\begin{pgfscope}%
\pgfsetbuttcap%
\pgfsetmiterjoin%
\definecolor{currentfill}{rgb}{1.000000,1.000000,1.000000}%
\pgfsetfillcolor{currentfill}%
\pgfsetlinewidth{0.000000pt}%
\definecolor{currentstroke}{rgb}{0.000000,0.000000,0.000000}%
\pgfsetstrokecolor{currentstroke}%
\pgfsetstrokeopacity{0.000000}%
\pgfsetdash{}{0pt}%
\pgfpathmoveto{\pgfqpoint{0.653890in}{0.322083in}}%
\pgfpathlineto{\pgfqpoint{3.753890in}{0.322083in}}%
\pgfpathlineto{\pgfqpoint{3.753890in}{3.402083in}}%
\pgfpathlineto{\pgfqpoint{0.653890in}{3.402083in}}%
\pgfpathlineto{\pgfqpoint{0.653890in}{0.322083in}}%
\pgfpathclose%
\pgfusepath{fill}%
\end{pgfscope}%
\begin{pgfscope}%
\pgfpathrectangle{\pgfqpoint{0.653890in}{0.322083in}}{\pgfqpoint{3.100000in}{3.080000in}}%
\pgfusepath{clip}%
\pgfsetbuttcap%
\pgfsetroundjoin%
\pgfsetlinewidth{0.501875pt}%
\definecolor{currentstroke}{rgb}{0.501961,0.501961,0.501961}%
\pgfsetstrokecolor{currentstroke}%
\pgfsetstrokeopacity{0.500000}%
\pgfsetdash{{1.850000pt}{0.800000pt}}{0.000000pt}%
\pgfpathmoveto{\pgfqpoint{2.203890in}{0.322083in}}%
\pgfpathlineto{\pgfqpoint{2.203890in}{3.402083in}}%
\pgfusepath{stroke}%
\end{pgfscope}%
\begin{pgfscope}%
\pgfsetbuttcap%
\pgfsetroundjoin%
\definecolor{currentfill}{rgb}{0.000000,0.000000,0.000000}%
\pgfsetfillcolor{currentfill}%
\pgfsetlinewidth{0.803000pt}%
\definecolor{currentstroke}{rgb}{0.000000,0.000000,0.000000}%
\pgfsetstrokecolor{currentstroke}%
\pgfsetdash{}{0pt}%
\pgfsys@defobject{currentmarker}{\pgfqpoint{0.000000in}{-0.048611in}}{\pgfqpoint{0.000000in}{0.000000in}}{%
\pgfpathmoveto{\pgfqpoint{0.000000in}{0.000000in}}%
\pgfpathlineto{\pgfqpoint{0.000000in}{-0.048611in}}%
\pgfusepath{stroke,fill}%
}%
\begin{pgfscope}%
\pgfsys@transformshift{2.203890in}{0.322083in}%
\pgfsys@useobject{currentmarker}{}%
\end{pgfscope}%
\end{pgfscope}%
\begin{pgfscope}%
\definecolor{textcolor}{rgb}{0.000000,0.000000,0.000000}%
\pgfsetstrokecolor{textcolor}%
\pgfsetfillcolor{textcolor}%
\pgftext[x=2.203890in,y=0.224861in,,top]{\color{textcolor}{\rmfamily\fontsize{10.000000}{12.000000}\selectfont\catcode`\^=\active\def^{\ifmmode\sp\else\^{}\fi}\catcode`\%=\active\def%{\%}Salarios Empleados}}%
\end{pgfscope}%
\begin{pgfscope}%
\pgfpathrectangle{\pgfqpoint{0.653890in}{0.322083in}}{\pgfqpoint{3.100000in}{3.080000in}}%
\pgfusepath{clip}%
\pgfsetbuttcap%
\pgfsetroundjoin%
\pgfsetlinewidth{0.501875pt}%
\definecolor{currentstroke}{rgb}{0.501961,0.501961,0.501961}%
\pgfsetstrokecolor{currentstroke}%
\pgfsetstrokeopacity{0.500000}%
\pgfsetdash{{1.850000pt}{0.800000pt}}{0.000000pt}%
\pgfpathmoveto{\pgfqpoint{0.653890in}{0.462083in}}%
\pgfpathlineto{\pgfqpoint{3.753890in}{0.462083in}}%
\pgfusepath{stroke}%
\end{pgfscope}%
\begin{pgfscope}%
\pgfsetbuttcap%
\pgfsetroundjoin%
\definecolor{currentfill}{rgb}{0.000000,0.000000,0.000000}%
\pgfsetfillcolor{currentfill}%
\pgfsetlinewidth{0.803000pt}%
\definecolor{currentstroke}{rgb}{0.000000,0.000000,0.000000}%
\pgfsetstrokecolor{currentstroke}%
\pgfsetdash{}{0pt}%
\pgfsys@defobject{currentmarker}{\pgfqpoint{-0.048611in}{0.000000in}}{\pgfqpoint{-0.000000in}{0.000000in}}{%
\pgfpathmoveto{\pgfqpoint{-0.000000in}{0.000000in}}%
\pgfpathlineto{\pgfqpoint{-0.048611in}{0.000000in}}%
\pgfusepath{stroke,fill}%
}%
\begin{pgfscope}%
\pgfsys@transformshift{0.653890in}{0.462083in}%
\pgfsys@useobject{currentmarker}{}%
\end{pgfscope}%
\end{pgfscope}%
\begin{pgfscope}%
\definecolor{textcolor}{rgb}{0.000000,0.000000,0.000000}%
\pgfsetstrokecolor{textcolor}%
\pgfsetfillcolor{textcolor}%
\pgftext[x=0.348333in, y=0.413889in, left, base]{\color{textcolor}{\rmfamily\fontsize{10.000000}{12.000000}\selectfont\catcode`\^=\active\def^{\ifmmode\sp\else\^{}\fi}\catcode`\%=\active\def%{\%}$\mathdefault{800}$}}%
\end{pgfscope}%
\begin{pgfscope}%
\pgfpathrectangle{\pgfqpoint{0.653890in}{0.322083in}}{\pgfqpoint{3.100000in}{3.080000in}}%
\pgfusepath{clip}%
\pgfsetbuttcap%
\pgfsetroundjoin%
\pgfsetlinewidth{0.501875pt}%
\definecolor{currentstroke}{rgb}{0.501961,0.501961,0.501961}%
\pgfsetstrokecolor{currentstroke}%
\pgfsetstrokeopacity{0.500000}%
\pgfsetdash{{1.850000pt}{0.800000pt}}{0.000000pt}%
\pgfpathmoveto{\pgfqpoint{0.653890in}{0.928750in}}%
\pgfpathlineto{\pgfqpoint{3.753890in}{0.928750in}}%
\pgfusepath{stroke}%
\end{pgfscope}%
\begin{pgfscope}%
\pgfsetbuttcap%
\pgfsetroundjoin%
\definecolor{currentfill}{rgb}{0.000000,0.000000,0.000000}%
\pgfsetfillcolor{currentfill}%
\pgfsetlinewidth{0.803000pt}%
\definecolor{currentstroke}{rgb}{0.000000,0.000000,0.000000}%
\pgfsetstrokecolor{currentstroke}%
\pgfsetdash{}{0pt}%
\pgfsys@defobject{currentmarker}{\pgfqpoint{-0.048611in}{0.000000in}}{\pgfqpoint{-0.000000in}{0.000000in}}{%
\pgfpathmoveto{\pgfqpoint{-0.000000in}{0.000000in}}%
\pgfpathlineto{\pgfqpoint{-0.048611in}{0.000000in}}%
\pgfusepath{stroke,fill}%
}%
\begin{pgfscope}%
\pgfsys@transformshift{0.653890in}{0.928750in}%
\pgfsys@useobject{currentmarker}{}%
\end{pgfscope}%
\end{pgfscope}%
\begin{pgfscope}%
\definecolor{textcolor}{rgb}{0.000000,0.000000,0.000000}%
\pgfsetstrokecolor{textcolor}%
\pgfsetfillcolor{textcolor}%
\pgftext[x=0.278889in, y=0.880555in, left, base]{\color{textcolor}{\rmfamily\fontsize{10.000000}{12.000000}\selectfont\catcode`\^=\active\def^{\ifmmode\sp\else\^{}\fi}\catcode`\%=\active\def%{\%}$\mathdefault{1000}$}}%
\end{pgfscope}%
\begin{pgfscope}%
\pgfpathrectangle{\pgfqpoint{0.653890in}{0.322083in}}{\pgfqpoint{3.100000in}{3.080000in}}%
\pgfusepath{clip}%
\pgfsetbuttcap%
\pgfsetroundjoin%
\pgfsetlinewidth{0.501875pt}%
\definecolor{currentstroke}{rgb}{0.501961,0.501961,0.501961}%
\pgfsetstrokecolor{currentstroke}%
\pgfsetstrokeopacity{0.500000}%
\pgfsetdash{{1.850000pt}{0.800000pt}}{0.000000pt}%
\pgfpathmoveto{\pgfqpoint{0.653890in}{1.395417in}}%
\pgfpathlineto{\pgfqpoint{3.753890in}{1.395417in}}%
\pgfusepath{stroke}%
\end{pgfscope}%
\begin{pgfscope}%
\pgfsetbuttcap%
\pgfsetroundjoin%
\definecolor{currentfill}{rgb}{0.000000,0.000000,0.000000}%
\pgfsetfillcolor{currentfill}%
\pgfsetlinewidth{0.803000pt}%
\definecolor{currentstroke}{rgb}{0.000000,0.000000,0.000000}%
\pgfsetstrokecolor{currentstroke}%
\pgfsetdash{}{0pt}%
\pgfsys@defobject{currentmarker}{\pgfqpoint{-0.048611in}{0.000000in}}{\pgfqpoint{-0.000000in}{0.000000in}}{%
\pgfpathmoveto{\pgfqpoint{-0.000000in}{0.000000in}}%
\pgfpathlineto{\pgfqpoint{-0.048611in}{0.000000in}}%
\pgfusepath{stroke,fill}%
}%
\begin{pgfscope}%
\pgfsys@transformshift{0.653890in}{1.395417in}%
\pgfsys@useobject{currentmarker}{}%
\end{pgfscope}%
\end{pgfscope}%
\begin{pgfscope}%
\definecolor{textcolor}{rgb}{0.000000,0.000000,0.000000}%
\pgfsetstrokecolor{textcolor}%
\pgfsetfillcolor{textcolor}%
\pgftext[x=0.278889in, y=1.347222in, left, base]{\color{textcolor}{\rmfamily\fontsize{10.000000}{12.000000}\selectfont\catcode`\^=\active\def^{\ifmmode\sp\else\^{}\fi}\catcode`\%=\active\def%{\%}$\mathdefault{1200}$}}%
\end{pgfscope}%
\begin{pgfscope}%
\pgfpathrectangle{\pgfqpoint{0.653890in}{0.322083in}}{\pgfqpoint{3.100000in}{3.080000in}}%
\pgfusepath{clip}%
\pgfsetbuttcap%
\pgfsetroundjoin%
\pgfsetlinewidth{0.501875pt}%
\definecolor{currentstroke}{rgb}{0.501961,0.501961,0.501961}%
\pgfsetstrokecolor{currentstroke}%
\pgfsetstrokeopacity{0.500000}%
\pgfsetdash{{1.850000pt}{0.800000pt}}{0.000000pt}%
\pgfpathmoveto{\pgfqpoint{0.653890in}{1.862083in}}%
\pgfpathlineto{\pgfqpoint{3.753890in}{1.862083in}}%
\pgfusepath{stroke}%
\end{pgfscope}%
\begin{pgfscope}%
\pgfsetbuttcap%
\pgfsetroundjoin%
\definecolor{currentfill}{rgb}{0.000000,0.000000,0.000000}%
\pgfsetfillcolor{currentfill}%
\pgfsetlinewidth{0.803000pt}%
\definecolor{currentstroke}{rgb}{0.000000,0.000000,0.000000}%
\pgfsetstrokecolor{currentstroke}%
\pgfsetdash{}{0pt}%
\pgfsys@defobject{currentmarker}{\pgfqpoint{-0.048611in}{0.000000in}}{\pgfqpoint{-0.000000in}{0.000000in}}{%
\pgfpathmoveto{\pgfqpoint{-0.000000in}{0.000000in}}%
\pgfpathlineto{\pgfqpoint{-0.048611in}{0.000000in}}%
\pgfusepath{stroke,fill}%
}%
\begin{pgfscope}%
\pgfsys@transformshift{0.653890in}{1.862083in}%
\pgfsys@useobject{currentmarker}{}%
\end{pgfscope}%
\end{pgfscope}%
\begin{pgfscope}%
\definecolor{textcolor}{rgb}{0.000000,0.000000,0.000000}%
\pgfsetstrokecolor{textcolor}%
\pgfsetfillcolor{textcolor}%
\pgftext[x=0.278889in, y=1.813889in, left, base]{\color{textcolor}{\rmfamily\fontsize{10.000000}{12.000000}\selectfont\catcode`\^=\active\def^{\ifmmode\sp\else\^{}\fi}\catcode`\%=\active\def%{\%}$\mathdefault{1400}$}}%
\end{pgfscope}%
\begin{pgfscope}%
\pgfpathrectangle{\pgfqpoint{0.653890in}{0.322083in}}{\pgfqpoint{3.100000in}{3.080000in}}%
\pgfusepath{clip}%
\pgfsetbuttcap%
\pgfsetroundjoin%
\pgfsetlinewidth{0.501875pt}%
\definecolor{currentstroke}{rgb}{0.501961,0.501961,0.501961}%
\pgfsetstrokecolor{currentstroke}%
\pgfsetstrokeopacity{0.500000}%
\pgfsetdash{{1.850000pt}{0.800000pt}}{0.000000pt}%
\pgfpathmoveto{\pgfqpoint{0.653890in}{2.328750in}}%
\pgfpathlineto{\pgfqpoint{3.753890in}{2.328750in}}%
\pgfusepath{stroke}%
\end{pgfscope}%
\begin{pgfscope}%
\pgfsetbuttcap%
\pgfsetroundjoin%
\definecolor{currentfill}{rgb}{0.000000,0.000000,0.000000}%
\pgfsetfillcolor{currentfill}%
\pgfsetlinewidth{0.803000pt}%
\definecolor{currentstroke}{rgb}{0.000000,0.000000,0.000000}%
\pgfsetstrokecolor{currentstroke}%
\pgfsetdash{}{0pt}%
\pgfsys@defobject{currentmarker}{\pgfqpoint{-0.048611in}{0.000000in}}{\pgfqpoint{-0.000000in}{0.000000in}}{%
\pgfpathmoveto{\pgfqpoint{-0.000000in}{0.000000in}}%
\pgfpathlineto{\pgfqpoint{-0.048611in}{0.000000in}}%
\pgfusepath{stroke,fill}%
}%
\begin{pgfscope}%
\pgfsys@transformshift{0.653890in}{2.328750in}%
\pgfsys@useobject{currentmarker}{}%
\end{pgfscope}%
\end{pgfscope}%
\begin{pgfscope}%
\definecolor{textcolor}{rgb}{0.000000,0.000000,0.000000}%
\pgfsetstrokecolor{textcolor}%
\pgfsetfillcolor{textcolor}%
\pgftext[x=0.278889in, y=2.280555in, left, base]{\color{textcolor}{\rmfamily\fontsize{10.000000}{12.000000}\selectfont\catcode`\^=\active\def^{\ifmmode\sp\else\^{}\fi}\catcode`\%=\active\def%{\%}$\mathdefault{1600}$}}%
\end{pgfscope}%
\begin{pgfscope}%
\pgfpathrectangle{\pgfqpoint{0.653890in}{0.322083in}}{\pgfqpoint{3.100000in}{3.080000in}}%
\pgfusepath{clip}%
\pgfsetbuttcap%
\pgfsetroundjoin%
\pgfsetlinewidth{0.501875pt}%
\definecolor{currentstroke}{rgb}{0.501961,0.501961,0.501961}%
\pgfsetstrokecolor{currentstroke}%
\pgfsetstrokeopacity{0.500000}%
\pgfsetdash{{1.850000pt}{0.800000pt}}{0.000000pt}%
\pgfpathmoveto{\pgfqpoint{0.653890in}{2.795417in}}%
\pgfpathlineto{\pgfqpoint{3.753890in}{2.795417in}}%
\pgfusepath{stroke}%
\end{pgfscope}%
\begin{pgfscope}%
\pgfsetbuttcap%
\pgfsetroundjoin%
\definecolor{currentfill}{rgb}{0.000000,0.000000,0.000000}%
\pgfsetfillcolor{currentfill}%
\pgfsetlinewidth{0.803000pt}%
\definecolor{currentstroke}{rgb}{0.000000,0.000000,0.000000}%
\pgfsetstrokecolor{currentstroke}%
\pgfsetdash{}{0pt}%
\pgfsys@defobject{currentmarker}{\pgfqpoint{-0.048611in}{0.000000in}}{\pgfqpoint{-0.000000in}{0.000000in}}{%
\pgfpathmoveto{\pgfqpoint{-0.000000in}{0.000000in}}%
\pgfpathlineto{\pgfqpoint{-0.048611in}{0.000000in}}%
\pgfusepath{stroke,fill}%
}%
\begin{pgfscope}%
\pgfsys@transformshift{0.653890in}{2.795417in}%
\pgfsys@useobject{currentmarker}{}%
\end{pgfscope}%
\end{pgfscope}%
\begin{pgfscope}%
\definecolor{textcolor}{rgb}{0.000000,0.000000,0.000000}%
\pgfsetstrokecolor{textcolor}%
\pgfsetfillcolor{textcolor}%
\pgftext[x=0.278889in, y=2.747222in, left, base]{\color{textcolor}{\rmfamily\fontsize{10.000000}{12.000000}\selectfont\catcode`\^=\active\def^{\ifmmode\sp\else\^{}\fi}\catcode`\%=\active\def%{\%}$\mathdefault{1800}$}}%
\end{pgfscope}%
\begin{pgfscope}%
\pgfpathrectangle{\pgfqpoint{0.653890in}{0.322083in}}{\pgfqpoint{3.100000in}{3.080000in}}%
\pgfusepath{clip}%
\pgfsetbuttcap%
\pgfsetroundjoin%
\pgfsetlinewidth{0.501875pt}%
\definecolor{currentstroke}{rgb}{0.501961,0.501961,0.501961}%
\pgfsetstrokecolor{currentstroke}%
\pgfsetstrokeopacity{0.500000}%
\pgfsetdash{{1.850000pt}{0.800000pt}}{0.000000pt}%
\pgfpathmoveto{\pgfqpoint{0.653890in}{3.262083in}}%
\pgfpathlineto{\pgfqpoint{3.753890in}{3.262083in}}%
\pgfusepath{stroke}%
\end{pgfscope}%
\begin{pgfscope}%
\pgfsetbuttcap%
\pgfsetroundjoin%
\definecolor{currentfill}{rgb}{0.000000,0.000000,0.000000}%
\pgfsetfillcolor{currentfill}%
\pgfsetlinewidth{0.803000pt}%
\definecolor{currentstroke}{rgb}{0.000000,0.000000,0.000000}%
\pgfsetstrokecolor{currentstroke}%
\pgfsetdash{}{0pt}%
\pgfsys@defobject{currentmarker}{\pgfqpoint{-0.048611in}{0.000000in}}{\pgfqpoint{-0.000000in}{0.000000in}}{%
\pgfpathmoveto{\pgfqpoint{-0.000000in}{0.000000in}}%
\pgfpathlineto{\pgfqpoint{-0.048611in}{0.000000in}}%
\pgfusepath{stroke,fill}%
}%
\begin{pgfscope}%
\pgfsys@transformshift{0.653890in}{3.262083in}%
\pgfsys@useobject{currentmarker}{}%
\end{pgfscope}%
\end{pgfscope}%
\begin{pgfscope}%
\definecolor{textcolor}{rgb}{0.000000,0.000000,0.000000}%
\pgfsetstrokecolor{textcolor}%
\pgfsetfillcolor{textcolor}%
\pgftext[x=0.278889in, y=3.213889in, left, base]{\color{textcolor}{\rmfamily\fontsize{10.000000}{12.000000}\selectfont\catcode`\^=\active\def^{\ifmmode\sp\else\^{}\fi}\catcode`\%=\active\def%{\%}$\mathdefault{2000}$}}%
\end{pgfscope}%
\begin{pgfscope}%
\definecolor{textcolor}{rgb}{0.000000,0.000000,0.000000}%
\pgfsetstrokecolor{textcolor}%
\pgfsetfillcolor{textcolor}%
\pgftext[x=0.223333in,y=1.862083in,,bottom,rotate=90.000000]{\color{textcolor}{\rmfamily\fontsize{10.000000}{12.000000}\selectfont\catcode`\^=\active\def^{\ifmmode\sp\else\^{}\fi}\catcode`\%=\active\def%{\%}Salario}}%
\end{pgfscope}%
\begin{pgfscope}%
\pgfpathrectangle{\pgfqpoint{0.653890in}{0.322083in}}{\pgfqpoint{3.100000in}{3.080000in}}%
\pgfusepath{clip}%
\pgfsetrectcap%
\pgfsetroundjoin%
\pgfsetlinewidth{1.003750pt}%
\definecolor{currentstroke}{rgb}{0.000000,0.000000,0.000000}%
\pgfsetstrokecolor{currentstroke}%
\pgfsetdash{}{0pt}%
\pgfpathmoveto{\pgfqpoint{1.971390in}{0.943501in}}%
\pgfpathlineto{\pgfqpoint{2.436390in}{0.943501in}}%
\pgfpathlineto{\pgfqpoint{2.436390in}{1.688660in}}%
\pgfpathlineto{\pgfqpoint{1.971390in}{1.688660in}}%
\pgfpathlineto{\pgfqpoint{1.971390in}{0.943501in}}%
\pgfusepath{stroke}%
\end{pgfscope}%
\begin{pgfscope}%
\pgfpathrectangle{\pgfqpoint{0.653890in}{0.322083in}}{\pgfqpoint{3.100000in}{3.080000in}}%
\pgfusepath{clip}%
\pgfsetrectcap%
\pgfsetroundjoin%
\pgfsetlinewidth{1.003750pt}%
\definecolor{currentstroke}{rgb}{0.000000,0.000000,0.000000}%
\pgfsetstrokecolor{currentstroke}%
\pgfsetdash{}{0pt}%
\pgfpathmoveto{\pgfqpoint{2.203890in}{0.943501in}}%
\pgfpathlineto{\pgfqpoint{2.203890in}{0.462083in}}%
\pgfusepath{stroke}%
\end{pgfscope}%
\begin{pgfscope}%
\pgfpathrectangle{\pgfqpoint{0.653890in}{0.322083in}}{\pgfqpoint{3.100000in}{3.080000in}}%
\pgfusepath{clip}%
\pgfsetrectcap%
\pgfsetroundjoin%
\pgfsetlinewidth{1.003750pt}%
\definecolor{currentstroke}{rgb}{0.000000,0.000000,0.000000}%
\pgfsetstrokecolor{currentstroke}%
\pgfsetdash{}{0pt}%
\pgfpathmoveto{\pgfqpoint{2.203890in}{1.688660in}}%
\pgfpathlineto{\pgfqpoint{2.203890in}{3.262083in}}%
\pgfusepath{stroke}%
\end{pgfscope}%
\begin{pgfscope}%
\pgfpathrectangle{\pgfqpoint{0.653890in}{0.322083in}}{\pgfqpoint{3.100000in}{3.080000in}}%
\pgfusepath{clip}%
\pgfsetrectcap%
\pgfsetroundjoin%
\pgfsetlinewidth{1.003750pt}%
\definecolor{currentstroke}{rgb}{0.000000,0.000000,0.000000}%
\pgfsetstrokecolor{currentstroke}%
\pgfsetdash{}{0pt}%
\pgfpathmoveto{\pgfqpoint{2.087640in}{0.462083in}}%
\pgfpathlineto{\pgfqpoint{2.320140in}{0.462083in}}%
\pgfusepath{stroke}%
\end{pgfscope}%
\begin{pgfscope}%
\pgfpathrectangle{\pgfqpoint{0.653890in}{0.322083in}}{\pgfqpoint{3.100000in}{3.080000in}}%
\pgfusepath{clip}%
\pgfsetrectcap%
\pgfsetroundjoin%
\pgfsetlinewidth{1.003750pt}%
\definecolor{currentstroke}{rgb}{0.000000,0.000000,0.000000}%
\pgfsetstrokecolor{currentstroke}%
\pgfsetdash{}{0pt}%
\pgfpathmoveto{\pgfqpoint{2.087640in}{3.262083in}}%
\pgfpathlineto{\pgfqpoint{2.320140in}{3.262083in}}%
\pgfusepath{stroke}%
\end{pgfscope}%
\begin{pgfscope}%
\pgfpathrectangle{\pgfqpoint{0.653890in}{0.322083in}}{\pgfqpoint{3.100000in}{3.080000in}}%
\pgfusepath{clip}%
\pgfsetbuttcap%
\pgfsetroundjoin%
\pgfsetlinewidth{1.505625pt}%
\definecolor{currentstroke}{rgb}{1.000000,0.000000,0.000000}%
\pgfsetstrokecolor{currentstroke}%
\pgfsetdash{{5.550000pt}{2.400000pt}}{0.000000pt}%
\pgfpathmoveto{\pgfqpoint{0.653890in}{0.698162in}}%
\pgfpathlineto{\pgfqpoint{3.753890in}{0.698162in}}%
\pgfusepath{stroke}%
\end{pgfscope}%
\begin{pgfscope}%
\pgfpathrectangle{\pgfqpoint{0.653890in}{0.322083in}}{\pgfqpoint{3.100000in}{3.080000in}}%
\pgfusepath{clip}%
\pgfsetbuttcap%
\pgfsetroundjoin%
\pgfsetlinewidth{1.505625pt}%
\definecolor{currentstroke}{rgb}{0.000000,0.501961,0.000000}%
\pgfsetstrokecolor{currentstroke}%
\pgfsetdash{{5.550000pt}{2.400000pt}}{0.000000pt}%
\pgfpathmoveto{\pgfqpoint{0.653890in}{0.943501in}}%
\pgfpathlineto{\pgfqpoint{3.753890in}{0.943501in}}%
\pgfusepath{stroke}%
\end{pgfscope}%
\begin{pgfscope}%
\pgfpathrectangle{\pgfqpoint{0.653890in}{0.322083in}}{\pgfqpoint{3.100000in}{3.080000in}}%
\pgfusepath{clip}%
\pgfsetbuttcap%
\pgfsetroundjoin%
\pgfsetlinewidth{1.003750pt}%
\definecolor{currentstroke}{rgb}{1.000000,0.498039,0.054902}%
\pgfsetstrokecolor{currentstroke}%
\pgfsetdash{}{0pt}%
\pgfpathmoveto{\pgfqpoint{1.971390in}{1.357432in}}%
\pgfpathlineto{\pgfqpoint{2.436390in}{1.357432in}}%
\pgfusepath{stroke}%
\end{pgfscope}%
\begin{pgfscope}%
\pgfsetrectcap%
\pgfsetmiterjoin%
\pgfsetlinewidth{0.803000pt}%
\definecolor{currentstroke}{rgb}{0.000000,0.000000,0.000000}%
\pgfsetstrokecolor{currentstroke}%
\pgfsetdash{}{0pt}%
\pgfpathmoveto{\pgfqpoint{0.653890in}{0.322083in}}%
\pgfpathlineto{\pgfqpoint{0.653890in}{3.402083in}}%
\pgfusepath{stroke}%
\end{pgfscope}%
\begin{pgfscope}%
\pgfsetrectcap%
\pgfsetmiterjoin%
\pgfsetlinewidth{0.803000pt}%
\definecolor{currentstroke}{rgb}{0.000000,0.000000,0.000000}%
\pgfsetstrokecolor{currentstroke}%
\pgfsetdash{}{0pt}%
\pgfpathmoveto{\pgfqpoint{3.753890in}{0.322083in}}%
\pgfpathlineto{\pgfqpoint{3.753890in}{3.402083in}}%
\pgfusepath{stroke}%
\end{pgfscope}%
\begin{pgfscope}%
\pgfsetrectcap%
\pgfsetmiterjoin%
\pgfsetlinewidth{0.803000pt}%
\definecolor{currentstroke}{rgb}{0.000000,0.000000,0.000000}%
\pgfsetstrokecolor{currentstroke}%
\pgfsetdash{}{0pt}%
\pgfpathmoveto{\pgfqpoint{0.653890in}{0.322083in}}%
\pgfpathlineto{\pgfqpoint{3.753890in}{0.322083in}}%
\pgfusepath{stroke}%
\end{pgfscope}%
\begin{pgfscope}%
\pgfsetrectcap%
\pgfsetmiterjoin%
\pgfsetlinewidth{0.803000pt}%
\definecolor{currentstroke}{rgb}{0.000000,0.000000,0.000000}%
\pgfsetstrokecolor{currentstroke}%
\pgfsetdash{}{0pt}%
\pgfpathmoveto{\pgfqpoint{0.653890in}{3.402083in}}%
\pgfpathlineto{\pgfqpoint{3.753890in}{3.402083in}}%
\pgfusepath{stroke}%
\end{pgfscope}%
\begin{pgfscope}%
\definecolor{textcolor}{rgb}{0.000000,0.000000,0.000000}%
\pgfsetstrokecolor{textcolor}%
\pgfsetfillcolor{textcolor}%
\pgftext[x=2.203890in,y=3.485417in,,base]{\color{textcolor}{\rmfamily\fontsize{12.000000}{14.400000}\selectfont\catcode`\^=\active\def^{\ifmmode\sp\else\^{}\fi}\catcode`\%=\active\def%{\%}Distribución Salarial basada en Intervalos}}%
\end{pgfscope}%
\begin{pgfscope}%
\pgfsetbuttcap%
\pgfsetmiterjoin%
\definecolor{currentfill}{rgb}{1.000000,1.000000,1.000000}%
\pgfsetfillcolor{currentfill}%
\pgfsetfillopacity{0.800000}%
\pgfsetlinewidth{1.003750pt}%
\definecolor{currentstroke}{rgb}{0.800000,0.800000,0.800000}%
\pgfsetstrokecolor{currentstroke}%
\pgfsetstrokeopacity{0.800000}%
\pgfsetdash{}{0pt}%
\pgfpathmoveto{\pgfqpoint{2.423473in}{2.903750in}}%
\pgfpathlineto{\pgfqpoint{3.656667in}{2.903750in}}%
\pgfpathquadraticcurveto{\pgfqpoint{3.684445in}{2.903750in}}{\pgfqpoint{3.684445in}{2.931528in}}%
\pgfpathlineto{\pgfqpoint{3.684445in}{3.304861in}}%
\pgfpathquadraticcurveto{\pgfqpoint{3.684445in}{3.332639in}}{\pgfqpoint{3.656667in}{3.332639in}}%
\pgfpathlineto{\pgfqpoint{2.423473in}{3.332639in}}%
\pgfpathquadraticcurveto{\pgfqpoint{2.395695in}{3.332639in}}{\pgfqpoint{2.395695in}{3.304861in}}%
\pgfpathlineto{\pgfqpoint{2.395695in}{2.931528in}}%
\pgfpathquadraticcurveto{\pgfqpoint{2.395695in}{2.903750in}}{\pgfqpoint{2.423473in}{2.903750in}}%
\pgfpathlineto{\pgfqpoint{2.423473in}{2.903750in}}%
\pgfpathclose%
\pgfusepath{stroke,fill}%
\end{pgfscope}%
\begin{pgfscope}%
\pgfsetbuttcap%
\pgfsetroundjoin%
\pgfsetlinewidth{1.505625pt}%
\definecolor{currentstroke}{rgb}{1.000000,0.000000,0.000000}%
\pgfsetstrokecolor{currentstroke}%
\pgfsetdash{{5.550000pt}{2.400000pt}}{0.000000pt}%
\pgfpathmoveto{\pgfqpoint{2.451251in}{3.228472in}}%
\pgfpathlineto{\pgfqpoint{2.590139in}{3.228472in}}%
\pgfpathlineto{\pgfqpoint{2.729028in}{3.228472in}}%
\pgfusepath{stroke}%
\end{pgfscope}%
\begin{pgfscope}%
\definecolor{textcolor}{rgb}{0.000000,0.000000,0.000000}%
\pgfsetstrokecolor{textcolor}%
\pgfsetfillcolor{textcolor}%
\pgftext[x=2.840139in,y=3.179861in,left,base]{\color{textcolor}{\rmfamily\fontsize{10.000000}{12.000000}\selectfont\catcode`\^=\active\def^{\ifmmode\sp\else\^{}\fi}\catcode`\%=\active\def%{\%}P12: 901.18}}%
\end{pgfscope}%
\begin{pgfscope}%
\pgfsetbuttcap%
\pgfsetroundjoin%
\pgfsetlinewidth{1.505625pt}%
\definecolor{currentstroke}{rgb}{0.000000,0.501961,0.000000}%
\pgfsetstrokecolor{currentstroke}%
\pgfsetdash{{5.550000pt}{2.400000pt}}{0.000000pt}%
\pgfpathmoveto{\pgfqpoint{2.451251in}{3.034861in}}%
\pgfpathlineto{\pgfqpoint{2.590139in}{3.034861in}}%
\pgfpathlineto{\pgfqpoint{2.729028in}{3.034861in}}%
\pgfusepath{stroke}%
\end{pgfscope}%
\begin{pgfscope}%
\definecolor{textcolor}{rgb}{0.000000,0.000000,0.000000}%
\pgfsetstrokecolor{textcolor}%
\pgfsetfillcolor{textcolor}%
\pgftext[x=2.840139in,y=2.986250in,left,base]{\color{textcolor}{\rmfamily\fontsize{10.000000}{12.000000}\selectfont\catcode`\^=\active\def^{\ifmmode\sp\else\^{}\fi}\catcode`\%=\active\def%{\%}P25: 1006.32}}%
\end{pgfscope}%
\end{pgfpicture}%
\makeatother%
\endgroup%

\end{figure}

Se nos cuestiono:

¿En caso de que se acuerde fijar el aumento en el salario tope de \mathdollar{1000}?

En ese caso abarcaria al \PorcentajeMenorAMil\% de los empleados.

\section{Coeficiente de Asimetria y Curtosis}

\newcommand{\CoeficienteFisherTipo}{Asimetria positiva}
\newcommand{\CoeficienteFisher}{0.98}
\newcommand{\CoeficientePearsonTipo}{Sesgo a la derecha}
\newcommand{\CoeficientePearson}{0.40}
\newcommand{\CoeficienteCurtosisTipo}{Leptocurtica}
\newcommand{\CoeficienteCurtosis}{3.43}

Se cuestiono si la distribucion de salarios es simétrica y que tan concentrado y/o dispersos se encuentran los datos (curtosis).

\subsection{Coeficientes de Asimetria}

Los valores obtenidos fueron:
\begin{itemize}
    \item Coeficiente de Fisher:
    \begin{itemize}
        \item Valor: \CoeficienteFisher
        \item Significado: \CoeficienteFisherTipo
    \end{itemize}
    \item Coeficiente de Pearson: 
    \begin{itemize}
        \item Valor: \CoeficientePearson
        \item Significado: \CoeficientePearsonTipo
    \end{itemize}
\end{itemize}

\subsection{Coeficiente de curtosis}

Se obtuvo que la curtosis es de \(\CoeficienteCurtosis\), es decir, es de tipo \CoeficienteCurtosisTipo

\emph{Ambos coeficientes se pueden observar en la figura \ref{fig:hist}}
\end{document}